\begin{examplebox}
	\textbf{\textcolor{CExample}{例1（基础）}}
	\textbf{\textcolor{CExample}{题目：}}
	求函数 $f(x)=2x^3-3x^2-12x+5$ 在闭区间 $[-2,3]$ 上的最大值和最小值。
	\textbf{\textcolor{CExample}{解题步骤：}}
	\begin{description}[style=nextline, leftmargin=2.8em, labelsep=0.5em, itemsep=0.45em]
		\item[\textbf{第一步（求可疑点）：}]
			求导得 $f'(x)=6x^2-6x-12=6(x-2)(x+1)$。令 $f'(x)=0$，解得 $x=-1$ 或 $x=2$，均在 $(-2,3)$ 内。没有不可导点，故可疑点为 $-2,-1,2,3$。
			\par\textbf{理由：}最值点只能在端点或导数为零的点中。
		\item[\textbf{第二步（计算函数值）：}]
			计算各可疑点函数值：
			\[
				f(-2)=1,\quad f(-1)=12,\quad f(2)=-15,\quad f(3)=-4.
			\]
			\par\textbf{理由：}代入函数表达式计算。
		\item[\textbf{第三步（比较得最值）：}]
			比较四个值：
			\[
				M_{\max}=12,\quad m_{\min}=-15.
			\]
			\par\textbf{理由：}最大值12在 $x=-1$ 处取得，最小值-15在 $x=2$ 处取得。
	\end{description}
	\textbf{\textcolor{CExample}{关键结论：}}
	\textbf{最大值 $12$，最小值 $-15$。}

	\medskip
	\textbf{\textcolor{CExample}{例2（稍变形）}}
	\textbf{\textcolor{CExample}{题目：}}
	求函数 $f(x)=|x^2-4|$ 在闭区间 $[-3,3]$ 上的最大值和最小值。
	\textbf{\textcolor{CExample}{解题步骤：}}
	\begin{description}[style=nextline, leftmargin=2.8em, labelsep=0.5em, itemsep=0.45em]
		\item[\textbf{第一步（消除绝对值）：}]
			写成分段函数：
			\[
				f(x)=
				\begin{cases}
					x^2-4, & |x|\ge2,\\ 4-x^2, & |x|<2.
			\end{cases}
			\]
			注意在 $x=\pm2$ 处 $f(x)$ 连续但不可导（尖点），所以它们也是可疑点。
			\par\textbf{理由：}绝对值函数在转折点处不可导。
		\item[\textbf{第二步（计算各点值）：}]
			计算端点、分段点以及不可导点处的函数值：
			\[
				f(-3)=5,\; f(3)=5,\; f(-2)=0,\; f(2)=0,\; f(0)=4.
			\]
			可疑点集合为 $\{-3,-2,0,2,3\}$。
			\par\textbf{理由：}端点和不可导点必须列入。
		\item[\textbf{第三步（比较得最值）：}]
			比较所有值得：
			\[
				M_{\max}=5,\quad m_{\min}=0.
			\]
			\par\textbf{理由：}最大值5在 $x=\pm3$ 处取得，最小值0在 $x=\pm2$ 处取得。
	\end{description}
	\textbf{\textcolor{CExample}{关键结论：}}
	\textbf{最大值 $5$，最小值 $0$。}

	\medskip
	\textbf{\textcolor{CExample}{例3（综合应用）}}
	\textbf{\textcolor{CExample}{题目：}}
	已知不等式 $x^3-3x+1\le a$ 对任意 $x\in[-1,2]$ 恒成立，求实数 $a$ 的最小值。
	\textbf{\textcolor{CExample}{解题步骤：}}
	\begin{description}[style=nextline, leftmargin=2.8em, labelsep=0.5em, itemsep=0.45em]
		\item[\textbf{第一步（构造函数）：}]
			令 $g(x)=x^3-3x+1$，则问题转化为 $a\ge \max_{x\in[-1,2]} g(x)$。
			\par\textbf{理由：}恒成立要求 $a$ 不小于 $g(x)$ 的最大值。
		\item[\textbf{第二步（求可疑点）：}]
			$g'(x)=3x^2-3=3(x-1)(x+1)$。在 $(-1,2)$ 内驻点为 $x=1$（$x=-1$ 是端点）。没有不可导点。可疑点集合为 $\{-1,1,2\}$。
			\par\textbf{理由：}最值在端点或驻点处取得。
		\item[\textbf{第三步（计算并比较）：}]
			计算各点值：
			\[
				g(-1)=3,\quad g(1)=-1,\quad g(2)=3.
			\]
			最大值为 $3$。
			\par\textbf{理由：}最大值为3。
		\item[\textbf{第四步（得参数范围）：}]
			由 $a\ge3$ 得 $a$ 的最小值为 $3$。
			\par\textbf{理由：}满足恒成立的最小 $a$ 等于最大值。
	\end{description}
	\textbf{\textcolor{CExample}{关键结论：}}
	\textbf{实数 $a$ 的最小值为 $3$。}
\end{examplebox}
